% ===== PRÉAMBULE CANONIQUE — à coller VERBATIM en tête de CHAQUE .tex =====
\documentclass[11pt,a4paper]{article}
\usepackage[utf8]{inputenc}\usepackage[T1]{fontenc}\usepackage{lmodern}
\usepackage[french]{babel}
\usepackage{amsmath,amssymb,mathtools}\usepackage{icomma}
\usepackage{siunitx}\sisetup{locale=FR}  % \qty{}{}, \si{}, \ang{}
\usepackage{xcolor}\usepackage{colortbl}
\usepackage{tikz}\usepackage{pgfplots}\pgfplotsset{compat=1.18}
\usepackage[siunitx,europeanresistors]{circuitikz} % circuits (élec.)
\usepackage[most]{tcolorbox}\usepackage{enumitem}\usepackage{array,booktabs}
\usepackage[a4paper,margin=2.2cm]{geometry}
\usetikzlibrary{arrows.meta,calc,angles,quotes,positioning,patterns,%
  backgrounds,shapes.geometric,decorations.pathmorphing,decorations.markings,intersections}
\definecolor{primary}{RGB}{222,49,99}\definecolor{primarydark}{RGB}{150,22,55}
\definecolor{defcol}{RGB}{0,114,178}\definecolor{methcol}{RGB}{52,92,148}
\definecolor{excol}{RGB}{0,135,90}\definecolor{attncol}{RGB}{200,40,55}
\tcbset{boxbase/.style={enhanced,breakable,boxrule=0.9pt,arc=2.5pt,
  left=3mm,right=3mm,top=2.3mm,bottom=2.3mm,fonttitle=\bfseries,coltitle=white}}
% --- boîtes TUTORIEL (noms EXACTS, ne pas renommer) ---
\newtcolorbox{cle}[1][]{boxbase,colback=defcol!6,colframe=defcol,
  title={\textsc{À retenir}\ifx\relax#1\relax\relax\else\textmd{ : \itshape #1}\fi}}
\newtcolorbox{methode}[1][]{boxbase,colback=methcol!7,colframe=methcol,sharp corners,
  title={\textsc{Méthode}\ifx\relax#1\relax\relax\else\textmd{ --- \itshape #1}\fi}}
\newtcolorbox{exemplebox}[1][]{boxbase,colback=excol!5,colframe=excol,
  title={\textsc{Exemple}\ifx\relax#1\relax\relax\else\textmd{ : \itshape #1}\fi}}
\newtcolorbox{piege}[1][]{boxbase,colback=attncol!6,colframe=attncol,
  title={\textsc{Piège}\ifx\relax#1\relax\relax\else\textmd{ : \itshape #1}\fi}}
\newtcolorbox{remarque}[1][]{boxbase,colback=black!4,colframe=black!45,
  title={\textsc{Remarque}\ifx\relax#1\relax\relax\else\textmd{ : \itshape #1}\fi}}
% --- boîtes EXERCICE / CORRIGÉ (noms EXACTS, auto-comptées) ---
\newtcolorbox[auto counter]{exo}[1][]{boxbase,colback=primary!4,colframe=primary,
  title={\textsc{Exercice}~\thetcbcounter\ifx\relax#1\relax\relax\else\textmd{ --- \itshape #1}\fi}}
\newtcolorbox[auto counter]{corrige}[1][]{boxbase,colback=excol!4,colframe=excol,
  title={\textsc{Corrigé}~\thetcbcounter\ifx\relax#1\relax\relax\else\textmd{ --- \itshape #1}\fi}}
\newcommand{\vv}[1]{\vec{#1}}\newcommand{\ds}{\displaystyle}
\newcommand{\R}{\mathbb{R}}\newcommand{\N}{\mathbb{N}}\newcommand{\Z}{\mathbb{Z}}
% (un \usepackage{fancyhdr}+en-tête est possible ; il est ignoré côté web)
% =========================================================================
\begin{document}
\begin{exo}[moment d'une force]
On visse un boulon avec une clé : une force $F=\qty{30}{N}$ est appliquée perpendiculairement à la
clé, à une distance $b=\qty{0.2}{m}$ de l'axe du boulon.
\begin{center}
\begin{tikzpicture}[>=Stealth,scale=1]
  \fill[black!70] (0,0) circle (0.12);
  \node[below=3pt] at (0,-0.1){\scriptsize axe};
  \draw[primary!70,line width=3pt] (0,0)--(2.6,0);
  \draw[<->,black!55] (0,-0.5)--(2.6,-0.5) node[midway,fill=white,inner sep=1pt]{\scriptsize $b$};
  \draw[->,line width=1.3pt,attncol] (2.6,0)--(2.6,1.1) node[above]{$\vv F$};
\end{tikzpicture}
\end{center}
\begin{enumerate}[label=\alph*)]
  \item Calcule le moment $\mathcal{M}$ de la force par rapport à l'axe.
  \item Dans quelle unité s'exprime-t-il ?
\end{enumerate}
\end{exo}

\begin{exo}[plus loin, plus efficace]
Sur la même clé, on applique toujours la même force $F=\qty{30}{N}$, mais à deux distances de l'axe :
$b_1=\qty{0.1}{m}$ puis $b_2=\qty{0.3}{m}$.
\begin{enumerate}[label=\alph*)]
  \item Calcule le moment dans chaque cas.
  \item Où faut-il pousser pour tourner le boulon le plus facilement ?
\end{enumerate}
\end{exo}

\begin{exo}[un moment nul]
Une force est appliquée sur une porte, mais sa \textbf{ligne d'action passe exactement par l'axe} des
gonds.
\begin{center}
\begin{tikzpicture}[>=Stealth,scale=1]
  \fill[black!70] (0,0) circle (0.1) node[left=3pt]{\scriptsize axe};
  \draw[primary!70,line width=2.5pt] (0,-0.05) rectangle (0.18,2);
  \draw[->,line width=1.3pt,attncol] (-1.3,1)--(-0.1,1) node[pos=0,left]{$\vv F$};
  \draw[black!45,densely dashed] (-1.4,1)--(0.4,1);
\end{tikzpicture}
\end{center}
\begin{enumerate}[label=\alph*)]
  \item Que vaut le bras de levier $b$ de cette force ?
  \item Que vaut son moment ? La porte tourne-t-elle ?
\end{enumerate}
\end{exo}

\begin{exo}[moment d'un couple]
Deux forces de même intensité $F=\qty{8}{N}$, parallèles et de sens opposé, forment un \textbf{couple}.
La distance entre leurs lignes d'action est $d=\qty{0.25}{m}$.
\begin{enumerate}[label=\alph*)]
  \item Que vaut la résultante (somme) des deux forces ?
  \item Calcule le moment du couple.
\end{enumerate}
\end{exo}

\begin{exo}[travail ou moment ?]
Deux grandeurs se calculent par « une force multipliée par une longueur » et se mesurent en
$\si{\newton}\times\si{\metre}$ : le \textbf{travail} d'une force et le \textbf{moment} d'une force.
\begin{enumerate}[label=\alph*)]
  \item Laquelle est une \emph{énergie}, exprimée en joules ?
  \item Pour laquelle la longueur utilisée est-elle \emph{parallèle} à la force ? Pour laquelle est-elle
        \emph{perpendiculaire} à la force ?
\end{enumerate}
\end{exo}

\end{document}
